\subsection{Stabilization and cutting} \label{subsection3.4}

In \cite{KO3} we discovered the stabilization and cut phenomena which are heavily used in our derivation of the set of def\/ining relations for the
diagonal reduction algebras of $\gl$-type. The consideration in \cite{KO3} uses the standard (by the f\/irst coordinates) embedding of $\gl_n$ into $\gl_{n+1}$.
In this subsection we shall make several more precise statements about the stabilization and cut considering now the embedding of $\gl_n\oplus\gl_1$
into $\gl_{n+1}$ (more generally, $\gl_n\oplus\gl_m$ into $\gl_{n+m}$). These precisions are needed to establish the behavior of the center of
the diagonal reduction algebra: namely we shall see that cutting preserves the centrality.


Notation: $\h$ in this subsection denotes the Cartan subalgebra of $\gl_{n+m}$.


Consider an embedding of $\gl_n\oplus \gl_m$ into $\gl_{n+m}$, given by an assignment $e_{ij}\mapsto e_{ij}$, $i,j=1,\ldots,n$, and
$e_{ab}\mapsto e_{n+a,n+b}$, $a,b=1,\ldots,m$, where $e_{kl}$ in the source are the generators of $\gl_n\oplus \gl_m$ and target $e_{kl}$
are in $\gl_{n+m}$. This rule together with the similar rule $E_{ij}\mapsto E_{ij}$ and $E_{ab}\mapsto E_{n+a,n+b}$ def\/ines an embedding of the Lie
algebra $(\gl_n\oplus \gl_m)\oplus(\gl_n\oplus \gl_m)$ into the Lie algebra $\gl_{n+m}\oplus\gl_{n+m}$ and of the enveloping algebras
$\Ar_n\otimes \Ar_m=\U(\gl_n\oplus\gl_n)\otimes \U(\gl_m\oplus\gl_m)$ into $\Ar_{n+m}=\U(\gl_{n+m}\oplus\gl_{n+m})$. This embedding clearly
maps nilpotent subalgebras of $\gl_n\oplus \gl_m$ to the corresponding nilpotent subalgebras of $\gl_{n+m}$ and thus def\/ines an embedding
$\iota_{n,m}:\Z_n\otimes\Z_m\to \Z_{n+m}$ of the corresponding double coset spaces. However, the map $\iota_{n,m}$ is not a~homomorphism of
algebras. This is because the multiplication maps are def\/ined with the help of projectors, which are dif\/ferent for $\gl_n\oplus \gl_m$ and $\gl_{n+m}$.


However, as we will explain now we can control certain dif\/ferences between the two multiplication maps. Let $\V_{n,m}$ be the left ideal of the
algebra $\Z_{n+m}$ generated by elements $\s_{ia}$ with $i=1,\ldots,n$ and $a=n+1,\ldots, n+m$; let  $\V_{n,m}'$ be the right ideal of the
algebra $\Z_{n+m}$ generated by elements $\s_{ai}$ with $i=1,\ldots, n$ and $a=n+1,\ldots, n+m$.


Write any element $\lambda\in Q_+$ (the positive cone of the root lattice of $\gl_{n+m}$) in the form $\lambda=\sum_{k=1}^{n+m}\lambda_k\varepsilon_k$.
The element $\lambda$ can be presented as a sum
\begin{gather}\label{lalaplapp}\lambda= \lambda'+\lambda''  ,\end{gather}
where $\lambda'$ is an element of the root lattice of $\gl_n\oplus \gl_m$, and $\lambda''$ is
proportional to the simple root $\varepsilon_n-\varepsilon_{n+1}$: $\lambda'=\sum_{k=1}^{n+m}\lambda'_k\varepsilon_k$ with
$\sum_{k=1}^n\lambda'_k =  \sum_{k=n+1}^{n+m}\lambda'_k=0$ and $\lambda''=c(\varepsilon_n-\varepsilon_{n+1})$.

\begin{lem}\label{lemma1}%\hspace{-.2cm}.
The left ideal  $\V_{n,m}\subset\Z_{n+m}$ consists of images in $\Z_{n+m}$ of  sums
$\sum_{ia}X_{ia} E_{ia}$ with $X_{ia}\in\Ab_{n+m}$, $i=1\lcd n$ and $a=n+1,\dots, n+m$.


The right ideal   $\V_{n,m}\subset\Z_{n+m}$ consists of images in $\Z_{n+m}$ of  sums $\sum_{ai} E_{ai}Y_{ai}$ with $Y_{ai}\in\Ab_{n+m}$,
$i=1\lcd n $ and  $a=n+1
\lcd n+m$.\end{lem}

\begin{proof}
Present the projector $P$ for the Lie algebra $\gl_{n+m}$ as a sum of terms
\begin{gather}
\label{wedeopro} \xi e_{-\gamma_1}\cdots e_{-\gamma_t}e_{\gamma'_{1}}\cdots e_{\gamma'_{t'}}  ,
\end{gather}
where $\xi\in\Uh$, $\gamma_1,\dots ,\gamma_t$ and $\gamma'_1,\dots ,\gamma'_{t'}$ are positive roots of $\gl_{n+m}$. For any $\lambda\in Q_+$
 denote by $P_\lambda$  the sum of above elements with $\gamma_1+\dots +\gamma_{t}=
\gamma'_1+\dots +\gamma'_{t'}=\lambda$. Then $P=\sum_{\lambda\in Q_+}P_\lambda$.
{} For any $X,Y\in\Ab$ def\/ine the element  $X\mult_\lambda Y$ as the image of $XP_\lambda Y$ in the reduction algebra. We have
$X\mult Y=\sum_{\lambda\in Q_+} X\mult_\lambda Y$.


For any $X\in\Ab_{n+m}$, $i=1\lcd n$ and $a=n+1\lcd n+m$ consider the product $X\mult_\lambda \s_{ia}$.


The product $X\mult_\lambda \s_{in}$ is zero if $\lambda''\not=0$ (the component $\lambda''$ is def\/ined by~(\ref{lalaplapp})).
Indeed, in this case in each summand of $P_\lambda$ one of $e_{\gamma'_{k'}}$ is equal to some $e_{jb}$, $j=1\lcd n$ and $b= n+1\lcd n+m$.
Choose an ordered basis of $\n_+$ which ends by all such $e_{jb}$ (ordered arbitrarily);
any element of~$\U(\n_+)$ can be written as a sum of ordered monomials, that is, monomials in which all such~$e_{jb}$ stand on the right.
Since $[e_{jb},E_{ia}]=0$ for any $i,j=1\lcd n$ and $a,b=n+1\lcd n+m$, the product
$e_{\gamma'_{k'}}E_{ia}$ belongs to the left ideal $\Ib$ and thus $X\mult_\lambda \s_{ia}=0$ in $\Z_{n+m}$.


If $\lambda''=0$ then generators of $\n_+$ in monomials entering the decomposition of  $P_\lambda$ are among the elements $e_{ij}$,
$1\leq i<j\leq n$, and $e_{ab}$, $n+1\leq a<b\leq n+m$ and thus their adjoint action leaves the space, spanned by all $E_{ia}$, $i=1\lcd n$,
$a=n+1\lcd n+m$ invariant, so $X\mult_\lambda \s_{ia}$ can be  presented as an image of the sum $\sum_{jb}X_{jb} E_{jb}$ with $X_{jb}\in\Ab_{n+m}$,
$j=1\lcd n$, $b=n+1\lcd n+m$. Thus, the left ideal, generated by all $\s_{ia}$ is contained in the vector space of images in $\Z_{n+m}$ of  sums
$\sum_{jb}X_{jb} E_{jb}$.


Moreover, for any $X\in\Ab_{n+m}$ the element $X\mult \s_{ia}$ is the image of $XE_{ia}+\sum_{j,b:\, j<i,\, b>a}\, X^{(jb)}E_{jb}$ for some $X^{(jb)}$
and the double induction on $i$ and $a$ proves the inverse inclusion.


The second part of lemma is proved similarly.
\end{proof}

\begin{cor}\label{corollary1}
We have the following decomposition of the free left $($and right$)$ $\Uh$-modules:
\begin{gather}\label{st4}
\Z_{n+m}= \mathrm{I}_{n,m}\oplus \Uh\cdot\iota_{n,m}(\Z_n\otimes\Z_m)   ,
\end{gather}
where $\mathrm{I}_{n,m}:=\V_{n,m}+\V'_{n,m}$.
\end{cor}

\begin{proof} The double coset space $\Z_{n+m}$ is a free left and right $\Uh$-module with a basis consisting of images of ordered monomials on elements
$E_{ij}$, $i,j=1\lcd n+m$; recall that we always use orders compatible with the partial order $<$ on $\h^*$, see (c) in Section~\ref{section2},
paragraph~2. We can choose an order for which all ordered monomials are of the form $XYZ$, where $X$ is a monomial on $E_{ai}$ with
$i=1\lcd n$ and $a=n+1\lcd n+m$, $Z$ is a monomial on $E_{ia}$  with $i=1\lcd n$ and $a=n+1\lcd n+m$ while $Y$ is a monomial on
$E_{ij}$ with $i,j=1\lcd n$ or $i,j=n+1\lcd n+m$. Then we apply the lemma above.
\end{proof}
